\section{Introduction}

Benchmarks of cultural knowledge report whether a model answers correctly.
They do not say how the answer is computed, or whether a correct answer about a festival in Bihar is produced the same way as one about a dish from Maharashtra.
Interpretability work on factual recall gives one way to ask this.
Reading the residual stream through the unembedding \citep{nostalgebraist2020logitlens} or through trained per-layer translators \citep{belrose2023tuned} shows the layer at which a model's prediction becomes stable, and causal methods \citep{meng2022locating,geva-etal-2023-dissecting} show which layers carry the information.

We use this to study Indian cultural knowledge with SANSKRITI \citep{maji-etal-2025-sanskriti}, a benchmark of 21{,}853 questions over \nStates{} states and union territories and sixteen cultural attributes.
For each correctly answered question we measure its \emph{crystallization depth}: the relative layer $d = \lstar / L$ at which the correct answer becomes, and remains, the model's choice.
The idea follows \citet{baldock2021difficulty}, who showed that the layer at which an image classifier's prediction settles tracks how hard or atypical an example is.

Our question is whether depth depends on how well documented the entity is.
Language models know less about rare entities \citep{kandpal2023longtail,mallen-etal-2023-trust}, and Indian regional culture is unevenly represented in web text.
If documentation matters beyond accuracy, then among questions a model gets right, answers about rare entities should form later, reflecting a longer computation or a weaker stored association.
If it does not, depth is a property of the model and the question form, and coverage gaps show up only as errors.
Both outcomes are informative for work on cultural representation in models \citep{yu2025cultrace,chelombitko2026cultural}, and we fixed the analysis, the primary metric and the criteria for a null result in advance (\S\ref{sec:prereg}).

Building the measurement exposed two problems that apply to any layer-wise study of multiple-choice answers.
First, in early and middle layers every model we tested ranks one particular answer letter first for nearly all questions.
A naive first-correct-layer metric therefore reports early crystallization exactly for questions whose gold answer happens to sit at that letter.
We remove this by presenting each question in its four cyclic option orders, so that each answer appears once at each letter, and averaging option probabilities by content.
Second, a single first-correct-layer index is noisy: in our pilot its test-retest correlation between two option orders was below 0.3.
We therefore use a continuous depth, the expected layer of the increase in gold probability, and choose between candidate metrics by split-half reliability under a rule written before the full run.

Our contributions are:
(1) a reliable, letter-bias-free measure of crystallization depth for multiple-choice questions, with its reliability reported;
(2) depth profiles for three model families on Indian cultural knowledge;
(3) a preregistered test of the link between entity documentation and depth, with controls for entity length and question type, and a state-level permutation test;
(4) an activation-patching check of whether lens depth matches the layers where the entity information is causally used; and
(5) a data-quality audit of SANSKRITI, including \nLeaks{} questions whose stem reveals the answer.

\todoresult{Final paragraph of the introduction, written after the run. If H1 is supported: state the size of the effect in layers per SD in each model and what it implies. If null: state the equivalence bound that was met and that depth is set by model and question form, not documentation. In either case, one sentence on patching.}
